\section{Study Design and Data}

The study focuses on Kazakhstan as a data-scarce urban setting in which city-level official totals are available but fine-resolution population labels for direct supervised modeling are not. This makes the setting scientifically useful for a bounded baseline paper: the task is not to recover a hidden census at grid level, but to construct a transparent proxy population surface whose city-level mass is anchored to official totals and whose intra-urban structure can be evaluated through multiple complementary evidence layers. The Kazakhstan framing is therefore part of the study design rather than only a geographic label. It provides a concrete national context in which open geospatial data, city-total calibration, and heterogeneous urban conditions can be examined together within a frozen and reproducible package.

In the frozen reference package, the city registry contains 10 cities with official totals available for calibrated inference: Almaty, Astana, Shymkent, Semey, Kurchatov, Taraz, Uralsk, Petropavlovsk, Ust Kamenogorsk, and Ridder. Of these, 8 belong to the validated baseline batch, meaning that they have the full combination of official totals, feature generation, and baseline-quality support required for the main v2 evidence stack. One city in the reference package, Semey, is additionally flagged as a smoke-passed example city, indicating that it was used for end-to-end runtime verification rather than as a separate scientific regime. The remaining cities in the registry still contribute to the frozen reference frame, but not every deeper validation layer is available uniformly across all 10 cities.

Official city totals are treated as calibration anchors throughout the manuscript. In the frozen workflow, these totals are compiled from the Bureau of National Statistics of Kazakhstan regional population sources, with a frozen reference date of 2026-02-01. Part of the registry is assembled from manually verified city records, and the remaining records are extracted from cached regional workbook sources within the same v2 pipeline. This provenance matters for the interpretation of the paper: the official totals are not incidental metadata, but the city-level accounting constraint around which the calibrated proxy surface is constructed. Supplementary Table~S2 summarizes the full frozen city registry, official-total provenance route, and completeness context retained with the manuscript package.

Features are derived from an OpenStreetMap-based stack spanning built-form intensity, floor-area proxies, transport structure, and point-of-interest or service-access variables. This design reflects the central premise of the paper: in the absence of true grid-level census labels, the goal is not direct truth reconstruction but a calibrated proxy population surface supported by transparent open-data inputs. The exact frozen feature inventory grouped by family is retained in Supplementary Table~S3.

The production configuration is fixed to a 500\,m grid and a random forest model, but the study design distinguishes between the full registry and the subset of cities used for deeper evidence layers. District benchmark comparison is available for Almaty, Astana, and Shymkent because those are the cities for which internal administrative reference material could be integrated into the v2 package. Qualitative validation is locked to Almaty and Astana, which together provide two deliberately contrasting reading contexts: Almaty as the stronger completeness case and Astana as the more cautious moderate-completeness case. These anchor-city choices are part of the paper's evidence design and should not be confused with claims that all cities are equally validated in the same way.

External structural comparison is conducted against WorldPop and GHS-POP. These products are used in the paper as independent gridded benchmarks rather than as cell-level truth sources. Their role in the study design is to provide an external frame of comparison for broad surface structure and hotspot-sensitive agreement, complementing the internal and qualitative layers without replacing them. In the same spirit, OSM completeness is carried as an interpretation context rather than as a training label or performance target. Stronger completeness supports more confident qualitative and internal readings, while weaker completeness motivates more cautious interpretation.

\begin{table}[H]
\centering
\caption{Study coverage summary for the frozen City1 v2 package. Counts are taken from the city-status registry, while completeness context is provided for the cities used in deeper validation layers.}
\label{tab:study-coverage}
\begin{tabular}{p{0.56\textwidth}p{0.34\textwidth}}
\toprule
Item & Value \\
\midrule
Cities in the frozen reference registry & 10 \\
Cities with official totals and calibrated inference support & 10 \\
Validated baseline cities & 8 \\
Smoke-passed example cities & 1 \\
District benchmark anchor cities & 3 (Almaty, Astana, and Shymkent) \\
Qualitative validation cities & 2 (Almaty and Astana) \\
OSM completeness context & Almaty: good (75.827); Astana: moderate (65.014); Shymkent: moderate (65.541) \\
\bottomrule
\end{tabular}
\end{table}

